\documentclass[a4paper, 11pt]{article}

% -- Coments
\usepackage{verbatim}

% ----- Fonts -----
% -- Fuente --
% \usepackage{fontspec}
% \setmonofont{JetBrainsMono Nerd Font}

\usepackage{babel}
\babelprovide[import, main]{spanish}

% -- Color --
\usepackage{xcolor}
%\definecolor{azul}{RGB}{00,33,99}
\definecolor{azul}{RGB}{35,72,180}

% -- Page Margin --
\usepackage[margin=1in]{geometry}

% -- Espaciados --
\newcommand{\Pspace}{0.5cm}
\newcommand{\Aspace}{0.2cm}

% -- Columnas --
\usepackage{multicol}

% -- Imagenes --
\usepackage{graphicx}
\usepackage{float}

% -- Matemáticas --
\usepackage{amsmath, amssymb}

% -- Gráficas --
\usepackage{pgfplots}
\pgfplotsset{compat=1.18}


\title
{
    Contexto Regional 2026-2 \\
    La Intendencia de Arizpe en la Independencia de Nueva España: 1810-1821 - Juan Domingo Vidargas del Moral \\

    Profesor José Moreno Vega
}

\begin{document}

\begin{otherlanguage}{spanish}
\maketitle
\end{otherlanguage}

\begin{center}
    \begin{tabular}{r|l}
        \textbf{Expediente} & \textbf{Nombre} \\ \hline
        219208106 & Bórquez Guerrero Angel Fernando \\
    \end{tabular}
\end{center}

{\Large 1. ¿Cuál es el argumento central del autor?} \\
Se explica que la guerra de independencia tuvo efectos particulares en la Intendencia de Arizpe y a diferencia de otras regiones de Nueva España, el movimiento insurgente no tuvo mucha presencia y organización en el noroeste. En un inicio, las autoridades de la intendencia se mantuvieron leales hacia la corona e incluso combatieron la insurgencia enviando dinero, tropas y apoyo para la defensa de San Blas. Sin embargo, la guerra también provocó cambios importantes en la región, como la alteración de las rutas comerciales y una mayor autonomía de las provincias internas. En 1821, varios funcionarios y grupos con poder económico se unieron al plan de iguala, ya que la nueva situación política podía favorecer sus intereses y permitirles ampliar su influencia regional.\par

\vspace{1cm}

{\Large 2. ¿Qué evidencias utiliza para respaldarlo?} \par
\begin{itemize}
    \item Documentos y registros de la época. Como las comunicaciones entre autoridades de la intendencia de Arizpe y los virreyes.
    \item Se presentan ejemplos de la participación de Arizpe en la guerra como el envío de dinero y hombres.
    \item Fechas de eventos como cuando se unieron al plan de Iguala.
\end{itemize}


\end{document}
